\title{Data Representation and Querying}
\subtitle{}
\author{ian.mcloughlin@gmit.ie}
\date{}


\begin{frame}
	\titlepage
\end{frame}

\begin{frame}
	\frametitle{Topics}
	\tableofcontents
\end{frame}

\section{About this module}

\begin{frame}{Learning outcomes}
	On completion of this module the learner will/should be able to:
	\begin{itemize}
		\item Explain the benefits and the limitations of a variety of data models.
		\item Determine the most appropriate data model given a set of requirements.
		\item Represent, integrate and querying large datasets using existing API's and frameworks.
		\item Describe the principles behind both the linked data and the open data movements.
	\end{itemize}
\end{frame}

\begin{frame}{Examinations}
	\begin{table}
		\begin{tabular}{p{4cm}r@{\hspace{0.5cm}}p{4cm}}
			Type & \% & Date \\
			\hline
			Project & 50 & Week 8 \\
			End of Semester Exam & 50 & See exams timetable
		\end{tabular}
	\end{table}
\end{frame}

\section{HTTP}
\begin{frame}{HyperText Transfer Protocol}
  \begin{description}
		\item[HyperText] Text with links.
    \vspace{0.25cm}
		\item[Transfer] Communication of data.
    \vspace{0.25cm}
		\item[Protocol] Set of rules for communication.
  \end{description}
\end{frame}

\begin{frame}{Uniform Resource Locator}
  \textcolor{amethyst}{http}://\textcolor{vermilion}{www}.\textcolor{azure(colorwheel)}{reddit.com}:\textcolor{brightpink}{80}\textcolor{green(ryb)}{/r/funny/}?\textcolor{harvardcrimson}{limit=1}
  
  \begin{table}
    \begin{tabular}{r@{\hspace{0.5cm}}p{6cm}}
      \textcolor{amethyst}{http} & Protocol \\
      \textcolor{brightpink}{80} & Port \\
      \textcolor{vermilion}{www} & Subdomain \\
      \textcolor{azure(colorwheel)}{reddit.com} & Domain \\
      \textcolor{green(ryb)}{/r/funny/} & Path \\
      \textcolor{harvardcrimson}{limit=1} & Parameter
    \end{tabular}
  \end{table}
\end{frame}


\begin{frame}{Resources}
 \begin{columns}[onlytextwidth]
   \begin{column}{0.5\textwidth}
     \centering
      \begin{figure}
        \begin{tikzpicture}
          \draw (0,0) -- (0,1.2) -- (0.7,1.2) -- (0.7,0.8) -- (1,0.8) -- (1,0) -- cycle;
          \draw (0.7,1.2) -- (1,0.8);
          \foreach \y in {0.2,0.4,0.6}{
            \draw (0.2,\y) -- (0.8,\y);
            \draw (0.2,0.8) -- (0.6,0.8);
            \draw (0.2,1) -- (0.6,1);
          }
        \end{tikzpicture}
        \caption*{File}
      \end{figure}
    \end{column}
    \begin{column}{0.5\textwidth}
      \centering
      \begin{figure}
      \begin{tikzpicture}
        \node (A) [cylinder, shape border rotate=90, draw,minimum height=1.3cm,minimum width=1cm] {};
      \end{tikzpicture}
      \caption*{Database}
      \end{figure}
    \end{column}
  \end{columns}
  
  \begin{quote}
    HTTP is used to transmit resources \ldots A resource is some chunk of information that can be identified by a URL \ldots The most common kind of resource is a file, but a resource may also be a dynamically-generated query result \ldots
  \end{quote}
  \citeurl{www.jmarshall.com/easy/http}
\end{frame}



\begin{frame}{HTTP Methods}
  \begin{description}
    \item[GET] Retrieve information from the server.
    \item[HEAD] Like get, but retrieve only the response header.
    \item[POST] Send data to the server.
    \item[PUT] Set the resource at the URI to the request data.
    \item[DELETE] Delete the resource at the URI.
    \item[CONNECT] Set up tunnel for other traffic to pass through HTTP.
    \item[OPTIONS] Find the allowable operations at the given URI.
    \item[TRACE] Echo the received request.
    \item[PATCH] Partial resource modification.
  \end{description}
\end{frame}


\begin{frame}[fragile]{Request (GET)}
  \begin{minted}{http}
GET /path/file.html HTTP/1.0
From: someuser@jmarshall.com
User-Agent: HTTPTool/1.0
  \end{minted}
\end{frame}

\begin{frame}[fragile]{Response}
  \begin{minted}{http}
HTTP/1.1 200 OK
Date: Mon, 27 Jul 2009 12:28:53 GMT
Server: Apache/2.2.14 (Win32)
Last-Modified: Wed, 22 Jul 2009 19:15:56 GMT
Content-Length: 88
Content-Type: text/html
Connection: Closed

<html>
  <body>
    <h1>Hello, World!</h1>
  </body>
</html>
  \end{minted}
\end{frame}


\begin{frame}[fragile]{Request (POST)}
  \begin{minted}{http}
POST /path/script.cgi HTTP/1.0
From: frog@jmarshall.com
User-Agent: HTTPTool/1.0
Content-Type: application/x-www-form-urlencoded
Content-Length: 32

home=Cosby&favorite+flavor=flies
  \end{minted}
\end{frame}

\begin{frame}{URL encoding}
  HTML form data is usually URL-encoded by changing;
  \begin{itemize}
    \item Unsafe characters to \% \emhp{xx} where \emph{xx} is the ascii value.
    \item All spaces to plusses.
    \item Names and values to: name1=value1\&name2=value2.
  \end{itemize}
  
  \begin{description}
    \item[GET] The parameters tag onto the URL with a ?, e.g.\ http://www.google.ie?q=Funny+cats.
    \item[POST] Parameters go in the body.
  \end{description}
  \citeurl{www.jmarshall.com/easy/http}
\end{frame}


\begin{frame}{Mindmap}
	\begin{adjustbox}{max totalsize={.9\textwidth}{.7\textheight},center}
		\begin{tikzpicture}[mindmap, grow cyclic, every node/.style=concept, concept color=orange!40,
		level 1/.append style={level distance=5cm,sibling angle=90},
		level 2/.append style={level distance=3cm,sibling angle=45}]
		
		\node{ShareLaTeX Tutorial Videos}
		child [concept color=blue!30] { node {Beginners Series}
			child { node {First Document}}
			child { node {Sections and Paragraphs}}
			child { node {Mathematics}}
			child { node {Images}}
			child { node {bibliography}}
			child { node {Tables and Matrices}}
			child { node {Longer Documents}}
		}
		child [concept color=yellow!30] { node {Thesis Series}
			child { node {Basic Structure}}
			child { node {Page Layout}}
			child { node {Figures, Subfigures and Tables}}
			child { node {Biblatex}}
			child { node {Title Page}}
		}
		child [concept color=teal!40] { node {Beamer Series}
			child { node {Getting Started}}
			child { node {Text, Pictures and Tables}}
			child { node {Blocks, Code and Hyperlinks}}
			child { node {Overlay Specifications}}
			child { node {Themes and Handouts}}
		}
		child [concept color=purple!50] { node {TikZ Series}
			child { node {Basic Drawing}}
			child { node {Geogebra}}
			child { node {Flow Charts}}
			child { node {Circuit Diagrams}}
			child [concept color=green!40] { node {Mind Maps}}
		};
		\end{tikzpicture}
	\end{adjustbox}
\end{frame}

\begin{frame}{Diagram}
	\tikzstyle{decision} = [diamond, draw, fill=blue!20, 
	text width=4.5em, text badly centered, node distance=3cm, inner sep=0pt]
	\tikzstyle{block} = [rectangle, draw, fill=blue!20, 
	text width=5em, text centered, rounded corners, minimum height=4em]
	\tikzstyle{line} = [draw, -latex']
	\tikzstyle{cloud} = [draw, ellipse,fill=red!20, node distance=3cm,
	minimum height=2em]
	
	\begin{adjustbox}{max totalsize={.9\textwidth}{.7\textheight},center} 
		\begin{tikzpicture}[node distance = 2cm, auto]
		% Place nodes
		\node [block] (init) {initialize model};
		\node [cloud, left of=init] (expert) {expert};
		\node [cloud, right of=init] (system) {system};
		\node [block, below of=init] (identify) {identify candidate models};
		\node [block, below of=identify] (evaluate) {evaluate candidate models};
		\node [block, left of=evaluate, node distance=3cm] (update) {update model};
		\node [decision, below of=evaluate] (decide) {is best candidate better?};
		\node [block, below of=decide, node distance=3cm] (stop) {stop};
		% Draw edges
		\path [line] (init) -- (identify);
		\path [line] (identify) -- (evaluate);
		\path [line] (evaluate) -- (decide);
		\path [line] (decide) -| node [near start] {yes} (update);
		\path [line] (update) |- (identify);
		\path [line] (decide) -- node {no}(stop);
		\path [line,dashed] (expert) -- (init);
		\path [line,dashed] (system) -- (init);
		\path [line,dashed] (system) |- (evaluate);
		\end{tikzpicture}
	\end{adjustbox}
\end{frame}